\chapter{How to read this book}\label{ch:howtoread}

\section*{What this book is}

This book turns two training courses into one mathematical textbook. The first course,
\emph{Finding anomalies in satellite imagery} (\file{docs/training/anomaly-detection/}, 34 parts),
teaches the domain, the classical detectors and the evaluation. The second,
\emph{Indradhanu, from zero} (\texttt{indradhanu workings/}, 71 parts), teaches the hyperspectral
foundation model the product uses to find anomalies, and the long history of experiments that
shaped it. Where the two courses covered the same ground, this book merges them into one chapter.

The courses already explain every idea in words and with small numbers. This book keeps that
and adds the mathematics underneath: each formula is derived one operation at a time, each rule
is checked on a tiny example the first time it is used, and each result is followed by a worked
example small enough to do on paper.

It is study material. It is subordinate to two things:
\begin{itemize}
  \item \textbf{The source code.} When this book and the code disagree, the code is right.
    Wherever a writer found a course and the code disagreeing, the chapter says so in a red
    \emph{Source check} box and follows the code.
  \item \textbf{The reference documentation}, \file{docs/01-orientation.md} through
    \file{docs/10-code-style.md}, and the low-level designs in \file{docs/lld/}.
\end{itemize}
The mathematics of naming a material from its spectrum (spectral library matching) has its own
companion volume in \file{research/math/spectral-matching/}; this book only points to it.

\section*{Who it is for}

You need school algebra: you can expand $(a+b)^2$, solve $2x + 3 = 11$, and read a sum written
with $\sum$. Everything else - vectors, matrices, eigenvalues, probability densities, derivatives
of functions of many variables, the parts of a neural network - is introduced from the beginning,
in the chapter that first needs it. If you learn slowly and carefully, the book is written for you.

\section*{How every chapter is built}

Every chapter has the same skeleton, so that you always know where you are.
\begin{enumerate}
  \item A green box, \textbf{The one thing this chapter teaches}, in one or two sentences.
  \item The motivation: what problem exists, and why the obvious fix fails.
  \item The mechanism, in numbered sections, with definitions before use, derivations one step
    per line, and purple \textbf{Worked example} boxes with every arithmetic step written out.
  \item \textbf{Common confusions}: the mistakes people actually make, as questions and answers.
  \item \textbf{Check yourself}: about five exercises, at least one by hand. Full solutions are in
    \cref{app:answers}.
\end{enumerate}
Other boxes you will meet:
\begin{itemize}
  \item a red \textbf{Source check} box, where a claim was checked against the code, or where a course
    and the code disagree;
  \item an \textbf{Analogy} box, one per chapter, tied to something outside software;
  \item an amber key-point box, the one sentence to remember from a section;
  \item a grey \textbf{Values that will drift} box, around numbers read from configurations,
    checkpoints and runs, with the date they were true.
\end{itemize}

\paragraph{Where the numbers come from.} A number about the project carries a small grey tag:
\tagM{file} means \emph{measured}, and names the file it was read from; \tagC{} means
\emph{computed} for this book, by hand in the text or by a script in \file{research/textbook/tools/};
\tagA{} means \emph{assumed} for illustration, with the assumption stated. Plain arithmetic inside a
worked example is not tagged.

\section*{Do the arithmetic}

The worked examples are deliberately tiny - three bands, a $2\times 2$ matrix, six tokens - so that
you can do them in a minute with a pencil. Reading them is not the same as doing them. Cover the
right-hand side of each line, compute it, then uncover it.

\section*{The parts, and what depends on what}

The book has ten parts. Part I is the domain: light, sensors and the data structure the product
hands to every algorithm. Part II is classical statistics: the linear algebra of covariance and
whitening, and the detectors built on it. Part III builds the Indradhanu network piece by piece,
with tensor shapes at every step and every parameter counted by hand; on the way it tours the five
convolutional autoencoders that came before it (\cref{ch:convae}), and it ends with Chakshu, the same machine
for one thermal band (\cref{ch:chakshu}). Part IV replaces Indradhanu's
spectral front end. Part V trains it. Part VI turns its output into an answer a user sees. Part VII
measures how good that answer is. Part VIII tells the history of the experiments as lessons, because
most of what the team knows was learned from runs that failed. Part IX covers what was built on 2026-10-31: the
spectral namer, a second model on the Indradhanu stem, and the first evaluation submission. Part X answers a question
the team kept being asked: why not use HyperSL, a published spectral foundation model with open weights? It shows what
HyperSL is and does well, then measures it on each of the product's jobs: library matching, anomaly detection, unmixing
and cloud masking.

\Cref{fig:howtoread:map} shows which chapters you need before which. A solid arrow means ``read
this first''. Read the book in order if you can; the map is for when you want to reach a chapter
sooner and need to know the shortest path to it.

\begin{figure}[htbp]
\centering
\begin{tikzpicture}[
  ch/.style={draw, rounded corners=2pt, minimum width=7mm, minimum height=5.5mm,
             font=\scriptsize\bfseries, inner sep=1pt},
  p1/.style={ch, fill=gray!12},
  p2/.style={ch, fill=loss!15},
  p3/.style={ch, fill=spat!12},
  p4/.style={ch, fill=spec!15},
  p5/.style={ch, fill=attn!12},
  p6/.style={ch, fill=loss!25},
  p7/.style={ch, fill=spec!25},
  p8/.style={ch, fill=gray!25},
  ar/.style={-{Stealth[length=1.6mm]}, gray!70!black, line width=0.5pt},
  lab/.style={font=\scriptsize\itshape, anchor=west, text=gray!40!black},
  x=9mm, y=8.5mm]
  % Part I
  \node[lab] at (-0.6,0.6) {I. Domain and data};
  \foreach \s/\x in {problem/0,light/1,sensors/2,product/3,bandpipeline/4} \node[p1] (\s) at (\x,0) {\ref{ch:\s}};
  \foreach \a/\b in {problem/light,light/sensors,sensors/product,product/bandpipeline} \draw[ar] (\a) -- (\b);
  % Part II
  \node[lab] at (5.4,0.6) {II. Statistical detection};
  \foreach \s/\x in {linalg/6,rx/7,localrx/8} \node[p2] (\s) at (\x,0) {\ref{ch:\s}};
  \draw[ar] (linalg) -- (rx); \draw[ar] (rx) -- (localrx);
  \draw[ar] (bandpipeline) -- (linalg);
  % Part III
  \node[lab] at (-0.6,-1.4) {III. Reconstruction, the Indradhanu architecture and Chakshu};
  \foreach \s/\x in {recon/0,shapes/1,convae/2,compressor/3,patchembed/4,attention/5,esa/6,block/7,encoder/8,masking/9,decoder/10,params/11,chakshu/12}
    \node[p3] (\s) at (\x,-2) {\ref{ch:\s}};
  \foreach \a/\b in {recon/shapes,shapes/convae,convae/compressor,compressor/patchembed,patchembed/attention,attention/esa,esa/block,block/encoder,encoder/masking,masking/decoder,decoder/params,params/chakshu}
    \draw[ar] (\a) -- (\b);
  \draw[ar] (bandpipeline) to[out=-120,in=60] (recon);
  \draw[ar] (linalg) to[out=-120,in=60] (shapes);
  % Part IV
  \node[lab] at (-0.6,-3.4) {IV. Spectral stem};
  \foreach \s/\x in {ceiling/0,crossstem/1,specdecoder/2} \node[p4] (\s) at (\x,-4) {\ref{ch:\s}};
  \draw[ar] (ceiling) -- (crossstem); \draw[ar] (crossstem) -- (specdecoder);
  \draw[ar] (params) to[out=-150,in=30] (ceiling);
  % Part V
  \node[lab] at (3.4,-3.4) {V. Training};
  \foreach \s/\x in {losses/4,trainloop/5,pretrain/6,seghead/7,finetune/8} \node[p5] (\s) at (\x,-4) {\ref{ch:\s}};
  \foreach \a/\b in {losses/trainloop,trainloop/pretrain,pretrain/seghead,seghead/finetune} \draw[ar] (\a) -- (\b);
  \draw[ar] (decoder) to[out=-120,in=60] (losses);
  \draw[ar] (specdecoder) to[out=-30,in=-150] (pretrain);
  % Part VI
  \node[lab] at (-0.6,-5.4) {VI. From model to answer};
  \foreach \s/\x in {inference/0,residualmaha/1,material/2} \node[p6] (\s) at (\x,-6) {\ref{ch:\s}};
  \draw[ar] (inference) -- (residualmaha); \draw[ar] (residualmaha) -- (material);
  \draw[ar] (trainloop) to[out=-150,in=30] (inference);
  \draw[ar] (localrx) to[out=-100,in=80] ([xshift=2mm]residualmaha.north);
  % Part VII
  \node[lab] at (3.4,-5.4) {VII. Evaluation};
  \foreach \s/\x in {measures/4,benchmarks/5,runeval/6,bootstrap/7} \node[p7] (\s) at (\x,-6) {\ref{ch:\s}};
  \foreach \a/\b in {measures/benchmarks,benchmarks/runeval,runeval/bootstrap} \draw[ar] (\a) -- (\b);
  \draw[ar] (residualmaha) to[out=-30,in=-150] (measures);
  % Part VIII
  \node[lab] at (-0.6,-7.4) {VIII. The experiments, told as lessons};
  \foreach \s/\x in {sortbug/0,ftruns/1,shipping/2,scoringgains/3,uncertainty/4,decoderwork/5,heavy/6}
    \node[p8] (\s) at (\x,-8) {\ref{ch:\s}};
  \foreach \a/\b in {sortbug/ftruns,ftruns/shipping,shipping/scoringgains,scoringgains/uncertainty,uncertainty/decoderwork,decoderwork/heavy}
    \draw[ar] (\a) -- (\b);
  \draw[ar] (bootstrap) to[out=-120,in=30] (sortbug);
  \draw[ar] (finetune) to[out=-60,in=60] (9,-6.5) to[out=-120,in=30] (ftruns);
  % Part IX
  \node[lab] at (7.5,-8.65) {IX. Naming and delivering};
  \node[p6] (namer) at (8,-8) {\ref{ch:namer}}; \node[p6] (submission) at (9,-8) {\ref{ch:submission}};
  \draw[ar] (namer) -- (submission);
  \draw[ar] (heavy) -- (namer);
  \draw[ar] (material) to[out=-15,in=175] (6.4,-6.75) to[out=-5,in=110] (namer);
  % Part X
  \node[lab] at (5.55,-10) {X. Why not HyperSL};
  \foreach \s/\x in {hslwhat/0,hslnorm/1,hsldetect/2,hslunmix/3,hslcloud/4,hslverdict/5}
    \node[p7] (\s) at (\x,-10) {\ref{ch:\s}};
  \foreach \a/\b in {hslwhat/hslnorm,hslnorm/hsldetect,hsldetect/hslunmix,hslunmix/hslcloud,hslcloud/hslverdict}
    \draw[ar] (\a) -- (\b);
  \draw[ar] (namer.south) to[out=-100,in=60] (hslwhat.north);
\end{tikzpicture}
\caption{The chapter dependency map. Each box is a chapter number; an arrow means ``read this
first''. Inside a part, read in order. The long arrows show the few places where a later part
leans on a specific earlier chapter: the linear algebra of \cref{ch:linalg} under the tensor
arithmetic of \cref{ch:shapes}, the local detectors of \cref{ch:localrx} under the residual
Mahalanobis score of \cref{ch:residualmaha}, and the fine-tuning machinery of \cref{ch:finetune}
under the experiment chapters. Part X leans on the namer (\cref{ch:namer}) and, through it, on the reconstruction,
masking, spectral-stem and evaluation chapters.}
\label{fig:howtoread:map}
\end{figure}

\section*{Five reading paths}

\begin{description}[style=nextline,leftmargin=1.2em,font=\normalfont\bfseries]
  \item[The whole book (recommended).] Chapters \ref{ch:problem} to \ref{ch:submission} in order. A fortnight of
    steady study.
  \item[Detection only.] Chapters \ref{ch:problem}--\ref{ch:localrx}, then \ref{ch:recon}, \ref{ch:residualmaha},
    \ref{ch:measures}--\ref{ch:bootstrap}. You learn the statistics of detection and how detections are measured,
    without the network.
  \item[The network only.] Chapters \ref{ch:problem}--\ref{ch:bandpipeline} for the data, \ref{ch:linalg} for the
    linear algebra, then \ref{ch:recon}--\ref{ch:finetune}.
  \item[Why not an outside foundation model.] Chapters \ref{ch:problem}--\ref{ch:bandpipeline}, \ref{ch:recon},
    \ref{ch:masking}, \ref{ch:specdecoder}, \ref{ch:residualmaha}, \ref{ch:material}, \ref{ch:measures} and
    \ref{ch:bootstrap}, then \ref{ch:hslwhat}--\ref{ch:hslverdict}. The case against HyperSL rests on why a pixel must
    not help rebuild itself, why statistics must be fixed by wavelength, and how detections are measured.
  \item[The thermal models.] Chapters \ref{ch:problem}--\ref{ch:bandpipeline}, then \ref{ch:recon}--\ref{ch:chakshu}:
    every architecture chapter, because Chakshu reuses them all; \ref{ch:convae} and \ref{ch:chakshu} are the thermal
    ones.
\end{description}

\section*{Numbers that will drift}

Architecture and mathematics age slowly. Run results do not: every pAUC, every validation loss,
every count of objects caught belongs to a run, a benchmark build and a date. Those numbers sit in
grey \emph{Values that will drift} boxes. The course indexes
(\file{docs/training/anomaly-detection/README.md} and \texttt{indradhanu workings/00-START-HERE.md})
list, part by part, which numbers to re-verify; the experiment READMEs under
\file{research/experiments/} are authoritative for every result.

\section*{The index and the lists}

After the table of contents come a \emph{List of worked examples} - use it to find an example to
redo when a later chapter assumes a calculation - and a list of figures. At the back are the full
solutions to every exercise (\cref{app:answers}), a glossary of every technical term with the
chapter that defines it (\cref{app:glossary}), the bibliography of the papers the courses cite, and
an index.
